% 固定六个有序输入，两组独立符号的教学实现，不改变真实标签。
% 阈值函数g_t在t左侧取-1、右侧取+1；两行逐组最大得分均为1/3。
\begin{tikzpicture}[foundation picture]
 \node[foundation heading] at (28,31) {随机符号};
 \node[foundation heading] at (86,31) {阈值空间};
 \node[foundation heading] at (144,31) {任意标签空间};
 \foreach \row/\case in {12/1,-42/2}{
  \node[foundation note,anchor=west] at (0,\row+18) {抽样\case};
  \foreach \offset in {0,58,116}{
   \begin{scope}[shift={(\offset,\row)}]
    \draw[foundation axis] (10,-14)--(47,-14) node[right] {$x_i$};
    \draw[foundation axis] (10,-14)--(10,11);
    \node[anchor=east] at (9,8) {$+1$};
    \node[anchor=east] at (9,-8) {$-1$};
    \foreach \i/\x in {1/14,2/20,3/26,4/32,5/38,6/44}{
     \node[below] at (\x,-14) {$\i$};
    }
   \end{scope}
  }
 }
 % 第一组符号为(+,+,-,-,+,+)：常数+1是一个最佳阈值函数。
 \draw[FigureRed,line width=1pt] (71,20)--(103,20);
 \draw[FigureRed,line width=1pt] (129,20)--(139,20)--(139,4)--(151,4)--(151,20)--(161,20);
 % 第二组符号为(-,+,-,+,-,+)：阈值在第三、四输入之间是一个最佳选择。
 \draw[FigureRed,line width=1pt] (71,-50)--(87,-50)--(87,-34)--(103,-34);
 \draw[FigureRed,line width=1pt] (129,-50)--(133,-50)--(133,-34)--(139,-34)--(139,-50)
  --(145,-50)--(145,-34)--(151,-34)--(151,-50)--(157,-50)--(157,-34)--(161,-34);
 \foreach \offset in {0,58,116}{
  \foreach \x/\s in {14/1,20/1,26/-1,32/-1,38/1,44/1}{
   \fill[FigureBlue,draw=FigureStroke,line width=.6pt] (\offset+\x,12+8*\s) circle (1.15);
  }
  \foreach \x/\s in {14/-1,20/1,26/-1,32/1,38/-1,44/1}{
   \fill[FigureBlue,draw=FigureStroke,line width=.6pt] (\offset+\x,-42+8*\s) circle (1.15);
  }
 }
 \foreach \x/\y in {84/4,90/4,78/-34,96/-50}{
  \draw[FigureYellow,line width=1.1pt] (\x,\y) circle (2.3);
 }
 \foreach \y in {-15,-69}{
  \node at (86,\y) {$\max_g\frac16\sum_i\sigma_i g(x_i)=\frac13$};
  \node at (144,\y) {$\max_g\frac16\sum_i\sigma_i g(x_i)=1$};
 }
\end{tikzpicture}%
